A foundation model-based agent operates in a perception-reasoning-action loop, using an LLM to process observations and execute actions (tool invocations, code execution, inter-agent communication). A multi-agent system (MAS) comprises multiple such agents, as exemplified by MetaGPT~\citep{metagpt2024}, CAMEL~\citep{camel2023}, ChatDev~\citep{chatdev2024}, and IoA~\citep{ioa2025}, operating within a shared environment or communicating to accomplish collective objectives~\citep{li2024, wang2026classical}. MAST~\citep{mast2025} identifies 14 recurring failure patterns in such systems. These failures motivate systematic diagnostics.

Multi-agent systems introduce three characteristics relevant to this survey~\citep{understandingmas2026}: (1)~\emph{distributed execution} across concurrent agents requiring coordination of LLM requests; (2)~\emph{inter-agent communication} creating complex traces that span agent boundaries, where failures cascade through communication channels~\citep{famas2025, traceelephant2026}; and (3)~\emph{inter-agent memory sharing}, where agents reuse prompts, tools, and cached state across the system. This last property, together with sparse activation, creates memory management challenges beyond single-model serving~\citep{zhou2026, kvcachesurvey2026}.

\paragraph*{Scope and verification procedure.} We focus on runtime operational layers distinct from three well-surveyed areas: orchestration~\citep{zhu2026orchestration}, single-model inference~\citep{kwon2023}, and offline evaluation~\citep{evalbenchmark2025}. Because we report numbers, we need an explicit rule for how far each number was checked rather than a closed list of trusted papers. We assign every cited system to a tier by the following procedure, applied in order:

\begin{enumerate}
\item \textbf{Tier 1a (full-text verified).} The source was read in full and the specific quantitative claim was located in the body text or a table. These numbers may be compared and quoted.
\item \textbf{Tier 1b (read in full, not page-verified).} The source was read in full and has a confirmed peer-reviewed venue, but the quantitative claim was taken from the published text without being cross-checked against a specific page, table, or figure. These numbers may be quoted with the $^\dagger$ marker but are not used to compute cross-system deltas.
\item \textbf{Tier 2 (abstract-only).} Only the abstract was read; the body was not examined. This applies to preprints and workshop papers, and equally to venue-published papers whose body we did not read. The claim is reported as stated and marked \emph{[a]}. The distinction from Tier 1b is therefore how much of the source was read, not the standing of its venue: Tier 1b means the full text was read, Tier 2 means it was not.
\item \textbf{Tier 3.} Documentation, platform pages, and vendor reports, labelled by source type and used for context only.
\end{enumerate}

Applying this procedure yields \textbf{8 Tier-1 papers across 9 table entries: 5 in Tier~1a} (FAMAS, TrajAudit, Agent Traces, TraceElephant, ScaleSim) and \textbf{3 in Tier~1b} (G-Safeguard at ACL 2025, ShieldAgent at ICML 2025, and AgenTracer at ICLR 2026). RootSE is the benchmark introduced by TrajAudit, so the two rows share one paper. Systems such as RAFFLES, AgentLocate, MultiAgentBench, and Who\&When do have confirmed venues but were read at abstract level, which under the procedure above places them in Tier~2, not Tier~1; the tier reflects verification depth rather than the standing of the venue. LlamaFirewall is likewise Tier~2: it is distributed as a preprint and was read at abstract level.